\chapter{완전한 설명은 어디까지 왔는가}
\label{chap:conclusion}

\section{현재의 결론}

2026년 8월 23일 현재, 자연어로 학습한 frontier-scale LLM 전체의 내부 알고리즘을 완전하고
충실하며 사람이 이해할 수 있는 형태로 발견한 연구는 없다. 현재 연구는 특정 행동의
부분 회로, 특정 prompt의 attribution graph, 일부 feature와 개입 가능한 방향을 제공하며,
논문들도 전역적 완전성을 주장하지 않는다.

반면 고정된 구현의 수치적 순전파는 완전히 명세할 수 있다. tokenizer, weights, input IDs,
dtype, kernel, random state와 decoding rule을 고정하면 embedding에서 모든 Q/K/V, residual,
MLP, logit과 token 선택까지 재현할 수 있다. 이 의미의 ``계산을 안다''와 학습된 의미
알고리즘을 안다는 주장을 혼동해서는 안 된다.

완전한 알고리즘을 아는 Transformer는 설계된 특수 사례에 존재한다. Tracr처럼 알려진
program을 compiler가 weight로 옮긴 모델은 source algorithm과 component 대응을 갖는다
\citep{lindner2023tracr}. 이 결과는 자연 학습 LLM을 완전히 해석했다는 결과가 아니라,
완전한 ground truth를 가진 실험실을 제공한다.

\section{완전성, 충실성, 압축성과 확장성}

유용한 완전한 설명은 적어도 네 축을 동시에 만족해야 한다. 한 축에서 좋은 결과를 다른
축의 증거로 대체할 수 없으므로, 각 축을 별도 metric과 실패 조건으로 관리해야 한다.

\begin{description}[style=nextline,leftmargin=2em]
  \item[완전성]
  target behavior에 필요한 계산을 빠뜨리지 않고 포함하는가? 선택한 분포 밖의 행동까지
  주장하려면 그 분포도 별도로 시험해야 한다.

  \item[인과적 충실성]
  설명 모델과 원 모델이 관측 출력뿐 아니라 관련 intervention과 counterfactual에 같은
  방향과 크기로 반응하는가?

  \item[압축성과 인간 이해 가능성]
  scalar operation의 전체 trace보다 짧으면서도 새로운 입력과 개입 결과를 예측할 수
  있는 변수, 절차와 불변량을 제공하는가?

  \item[확장성]
  소형 모델의 한 과제에 그치지 않고 더 큰 모델, 긴 context, 여러 언어와 행동에 적용할
  수 있으며 계산비용과 사람 검토 비용이 감당 가능한가?
\end{description}

이 네 축에는 tension이 존재한다. 모든 multiplication을 기록하면 완전성은 높지만 설명은
압축되지 않으며, head 몇 개에 짧은 이름을 붙이면 이해하기 쉽지만 누락된 feature와 경로가
많을 수 있다. replacement model은 해석 가능한 단위를 늘리지만 reconstruction error와
mechanistic mismatch를 도입할 수 있다 \citep{ameisen2025circuittracing}.

설명의 유일성도 보장되지 않는다. Méloux 등은 Boolean function과 소형 MLP의 후보 설명을
열거한 실험에서 같은 행동을 재현하는 여러 회로, 한 회로의 여러 해석과 같은 algorithm의
여러 subspace alignment가 존재하는 비식별성 사례를 보였다 \citep{meloux2025identifiable}.
따라서 ``유일한 진짜 이야기''보다 반사실 예측력, 조작 가능성, 간결성과 범위를 명시한
여러 동등 설명을 평가해야 할 수 있다.

\section{알려진 것의 경계}

현재 지식의 경계는 ``아는가, 모르는가''의 이분법보다 설명 수준별로 기록하는 편이 정확하다.
\Cref{tab:known-boundary}는 이 책의 결론을 단계별로 요약한다.

\begin{longtable}{@{}p{0.18\textwidth}p{0.34\textwidth}p{0.38\textwidth}@{}}
\caption{현대 decoder-only LLM에 대해 알려진 것과 남은 문제}
\label{tab:known-boundary}\\
\toprule
수준 & 강하게 알려진 내용 & 해결되지 않은 내용 \\
\midrule
\endfirsthead
\toprule
수준 & 강하게 알려진 내용 & 해결되지 않은 내용 \\
\midrule
\endhead
구현과 수학 & 공개 코드와 checkpoint의 순전파, tensor shape, logits와 decoding을 재현할 수 있다. & 비공개 training data와 recipe, hardware별 bit-wise 차이는 접근권한과 기록에 달려 있다. \\
표현 & 많은 변수는 특정 activation에서 probe되며 일부 방향은 steering할 수 있다. & 자연스러운 계산 단위의 유일한 dictionary와 전역 의미는 확립되지 않았다. \\
국소 인과 & ablation, patching과 editing으로 특정 input과 behavior의 구성요소 효과를 측정할 수 있다. & baseline, metric, self-repair와 distribution shift에 강건한 표준은 발전 중이다. \\
회로 & induction, IOI, factual recall과 여러 사례에서 구성요소 협력이 발견되었다. & 행동 전반을 덮는 완전하고 모델 간에 안정적인 회로 지도는 없다. \\
대형 모델 & feature 기반 attribution graph로 선택한 prompt의 중간 계산을 부분적으로 추적할 수 있다. & replacement error, attention QK, graph pruning과 성공 사례 편향이 남는다. \\
전역 알고리즘 & 설계·컴파일한 소형 Transformer에는 ground-truth program이 존재한다. & 자연 학습 frontier LLM 전체의 충실하고 압축된 algorithmic description은 없다. \\
\bottomrule
\end{longtable}

``모른다''는 문장은 아무것도 발견하지 못했다는 뜻이 아니다. induction head와 IOI,
subject enrichment와 attribute extraction, causal abstraction, SAE feature와 attribution
graph는 구체적인 진전이다. 다만 발견 범위를 보존해야 이 조각들을 잘못된 보편 이론으로
조기에 굳히지 않을 수 있다.

\section{연구 프로그램}

완전한 설명을 향한 실용적 연구 프로그램은 ground truth가 있는 작은 모델에서 평가 기준을
확립한 뒤 자연 모델과 더 큰 행동으로 확장해야 한다. 다음 여섯 단계는 이 책의 Phase 1부터
7까지를 연구 작업으로 변환한다.

\begin{enumerate}
  \item \textbf{실행을 고정한다.} 공개 model과 tokenizer의 순전파를 직접 구현하고 모든
        hook의 tensor를 reference code와 비교한다.
  \item \textbf{행동을 좁게 정의한다.} dataset, prompt generator, target token과 metric을
        versioned artifact로 만든다.
  \item \textbf{표현 가설을 만든다.} probe, geometry와 SAE로 후보 variable을 찾되, 이
        단계의 결론을 decodability로 제한한다.
  \item \textbf{경로 가설을 만든다.} QK/OV, DLA와 자동 검색으로 sender, receiver와
        intermediate variable을 명시한다.
  \item \textbf{반사실로 반증을 시도한다.} activation과 path intervention, causal
        abstraction, adversarial prompt와 model transfer로 예측을 시험한다.
  \item \textbf{설명을 압축하고 감사한다.} graph가 보존하는 행동, reconstruction error,
        누락 경로, 사람 해석 일치도와 계산비용을 함께 보고한다.
\end{enumerate}

각 단계의 산출물은 논문 문장보다 실행 가능한 기록이어야 한다. commit hash, checkpoint
digest, token IDs, hook name, seed, dtype, raw metric과 실패 사례를 보존하면 설명이 바뀔 때
어떤 근거가 수정되었는지 추적할 수 있다.

첫 번째 장기 project는 ``한 token을 완전히 추적하기''가 적합하다. 소형 공개 LLM의 한
factual prompt에서 embedding부터 target logit까지 모든 activation과 direct contribution을
저장하고, Geva의 세 단계가 paraphrase와 counterfactual에서 유지되는지 검증한다.

두 번째 장기 project는 ``알려진 program을 blind하게 재발견하기''가 적합하다. Tracr 모델의
source program을 가린 상태에서 probe, patching, ACDC와 SAE로 회로를 찾고, ground truth와
정밀도와 재현율을 비교한다. 이 project는 자연 LLM에서 정답이 없는 평가를 시작하기 전에
방법의 false positive와 false negative를 측정하게 해 준다.

\section{종료 조건이 아니라 갱신 규칙}

이 책은 ``LLM을 완전히 이해했다''는 종료 선언보다 근거가 들어올 때 설명을 고치는 갱신
규칙을 제공한다. 각 주장은 source, model, dataset, intervention, metric과 반증 조건을
가지며, 더 강한 결과가 나오면 claim ledger의 범위를 넓히거나 기존 문장을 철회한다.

새 논문은 headline이 아니라 증거 사슬로 편입해야 한다. 공개 구현과 data가 있는지,
baseline에 강건한지, 독립 model에서 재현되는지, replacement model의 오차가 포함되는지,
결과가 correlation인지 causation인지 확인한 뒤 해당 상자 등급을 정한다.

완전한 이해라는 목표는 현재의 불완전성을 숨기지 않을 때만 과학적 연구 계획이 된다.
우리가 완전히 아는 순전파부터 시작하여 표현, 국소 인과, 회로와 행동 일반화를 차례로
검증하면, 아직 존재하지 않는 전역 설명을 이미 발견한 것처럼 말하지 않으면서도 그 목표를
향해 누적적으로 전진할 수 있다.

\begin{factbox}
계산 자체는 재현 가능하고, 학습된 메커니즘은 부분적으로만 해석되었으며, 완전한 설명은
열린 연구 문제다. 앞으로의 진전은 더 많은 그럴듯한 이름이
아니라 더 강한 반사실 예측, 더 작은 설명 격차와 더 넓은 재현 범위로 측정해야 한다.
\end{factbox}
